\subsection{Discussion}
In this chapter, we first briefly described several tasks 
that we believe should be considered when it comes to AI for software engineering,
focusing on formal verification as one such task. 
We also highlighted a few cross-cutting challenges that permeate throughout many 
software-engineering tasks. 
Finally, to drive progress in AI for code, we pinpointed a small set of promising research directions for 
alleviating these challenges and advancing AI towards being a more capable software engineer.

In many ways, designing a roadmap and designing an evaluation are closely intertwined.
An evaluation makes a capability gap visible and concrete while giving the community a
clear way to measure progress.
A roadmap then synthesizes various evaluations into a 
broader view of where the field should go next, while at the same time 
identifying important capabilities that the community has not yet 
measured.
In this sense, evaluations and roadmaps provide a feedback loop for each other.
Evaluations reveal gaps in today's systems, roadmaps organize these gaps into 
a set of challenges and paths forward, which then motivate the next generation of
evaluations, and so on.

The roadmap in this chapter exemplifies this philosophy, as benchmarks discussed throughout this thesis played a key role in informing its construction.
For example, we built CRUXEval because existing benchmarks at the time did not capture
the natural programming capability of reasoning about code execution. Code execution and reasoning would 
have been a natural roadmap item then, as it was a simple yet important capability that was difficult for state-of-the-art models at the time. After CRUXEval, we saw in Section~\ref{sec:cruxeval-execution-impact} 
that the community began to develop new execution-guided techniques, suggesting 
broader directions that we discuss in our roadmap such as neurosymbolic methods and combining language
models with program-analysis tools.

More broadly, many of the challenges in our roadmap are difficult precisely because they
are not easily measurable. For measurable challenges like code execution or program repair,
it is relatively straightforward to design benchmarks that clearly show the gap, making 
it easier to make progress on the challenge and iterate toward better systems. In contrast, for
out-of-distribution domains like designing new programming languages or new program logics,
even defining the right benchmark can be part of the challenge. Overall, while some challenges are
known gaps that need better solutions, others need to begin with developing a convincing measurement.

The three lessons in this thesis help bridge benchmarking with roadmap design. 
First, choosing the right metric determines which missing capabilities become visible to the research community.
Second, looking beyond the score translates benchmark failures into actionable research directions and 
challenges.
Finally, understanding that what we measure shapes what we build reinforces the idea that evaluations and 
roadmaps are closely intertwined. 
Overall, benchmarks and analyses reveal where current systems fall short, while roadmaps 
help decide which gaps matter most and how the field might make progress toward closing them.
